\documentclass{article}
\usepackage{amsmath,amssymb}

\begin{document}

\begin{center}
  {\large\textbf{Hybridization vs.\ Co-Presence in Compositional Diffusion Models}}

  \medskip
  {\normalsize Author One$^1$, Author Two$^1$, Author Three$^2$}

  \smallskip
  {\small $^1$Institution One \quad $^2$Institution Two}

  \smallskip
  {\small Paper ID: XXXX}
\end{center}

\bigskip

Text-to-image diffusion models achieve \emph{semantic composition}
through natural-language descriptions, yet depend on training coverage
of specific concept pairings, and struggle with complex compositional
tasks such as counting or fine-grained attribute binding.
A parallel line of work pursues \emph{logical composition}: composing
smaller pretrained models at inference time, each specialising in a
distinct concept or constraint --- quick, data-efficient, and
generalisable to unseen combinations, yet seeking a compromise region
in the joint probability density that can blend rather than cleanly
separate concepts.
Both are used to compose complex scenes, yet one conditions generation
on a learned semantic meaning of the combined phrase while the other
composes their score functions directly.
We show they are not the same operation, and that conflating them
leads to wrong conclusions about what each paradigm can and cannot express.
Through trajectory-level analysis of the denoising flow, we demonstrate
that the two paradigms diverge early and consistently in latent space.
We characterise this gap by identifying the text prompt stable diffusion
would need to reproduce a given logical composition, then show that
linearly interpolating their velocities strikes a balance --- leveraging
the semantic coherence of language-guided generation and the
compositional expressibility of logical operators.

\bigskip
\noindent\textbf{Keywords:} compositional generation, diffusion models,
latent trajectories, approximate inversion, composability gap,
logical operators, score composition, flow matching

\end{document}
